\documentclass[11pt]{article}
\usepackage[letterpaper,margin=0.68in,headheight=15pt]{geometry}
\usepackage[T1]{fontenc}
\usepackage{lmodern}
\usepackage{amsmath,amssymb}
\usepackage{array,booktabs,enumitem,multicol}
\usepackage{xcolor}
\usepackage{tikz}
\usepackage{fancyhdr}
\usepackage{microtype}
\setlength{\parindent}{0pt}
\setlength{\parskip}{5pt}
\definecolor{olive}{HTML}{555B35}
\definecolor{gold}{HTML}{B18A22}
\definecolor{soft}{HTML}{F4F1DE}
\definecolor{bluegray}{HTML}{E8EEF2}
\pagestyle{fancy}
\fancyhf{}
\fancyhead[L]{\textcolor{olive}{\small MATH\& 141: Precalculus I}}
\fancyhead[R]{\textcolor{gold}{\small Fall 2026}}
\fancyfoot[C]{\thepage}
\renewcommand{\headrulewidth}{0.45pt}
\renewcommand{\footrulewidth}{0pt}
\setlist[enumerate]{leftmargin=*,itemsep=0.08in,topsep=0.05in}
\newcommand{\assessmenttitle}[2]{%
  {\Large\bfseries Module #1: #2}\par\vspace{3pt}
}
\newcommand{\studentline}{%
\noindent\textbf{Name:}\hspace{2.6in}\textbf{Date:}\hspace{1.15in}\par\vspace{7pt}}
\newcommand{\directions}[1]{\textit{#1}\par\vspace{4pt}}
\newcommand{\answerblank}[1]{\vspace{#1}}
\newcommand{\blankgrid}[4]{%
\begin{center}
\begin{tikzpicture}[x=0.75cm,y=0.75cm]
  \draw[step=1,gray!28,very thin] (#1,#3) grid (#2,#4);
  \draw[->,thick] (#1,0)--(#2,0) node[right] {$x$};
  \draw[->,thick] (0,#3)--(0,#4) node[above] {$y$};
  \foreach \x in {#1,...,#2}{\ifnum\x=0\else\draw (\x,0.10)--(\x,-0.10) node[below=2pt,font=\scriptsize]{\x};\fi}
  \foreach \y in {#3,...,#4}{\ifnum\y=0\else\draw (0.10,\y)--(-0.10,\y) node[left=2pt,font=\scriptsize]{\y};\fi}
\end{tikzpicture}
\end{center}}
\newcommand{\smallgrid}[4]{%
\begin{tikzpicture}[x=0.50cm,y=0.50cm]
  \draw[step=1,gray!28,very thin] (#1,#3) grid (#2,#4);
  \draw[->,thick] (#1,0)--(#2,0) node[right] {$x$};
  \draw[->,thick] (0,#3)--(0,#4) node[above] {$y$};
\end{tikzpicture}}
\begin{document}

% MODULE 5 PREQUIZ
\assessmenttitle{5}{Inverses and Radicals -- Prequiz}
\studentline
\directions{Four short questions. Remember that an inverse function reverses the original input-output rule.}
\begin{enumerate}
\item Find $f^{-1}(x)$ for $f(x)=4x+1$.
\answerblank{0.50in}
\item If $f(x)=2x-3$, find $f^{-1}(9)$.
\answerblank{0.48in}
\item Explain why $g(x)=x^2+4$ does not have an inverse function on all real numbers. If the domain is restricted to $x\ge 0$, find $g^{-1}(x)$ and state its domain.
\answerblank{0.82in}
\item Solve $\sqrt{3x+1}=7$ and check your solution.
\answerblank{0.70in}
\end{enumerate}
\newpage

% MODULE 5 CHECKIN PAGE 1
\assessmenttitle{5}{Inverses and Radicals -- Weekly Check-In}
\studentline
\directions{Three questions. Remember that an inverse reverses the original input-output rule.}
\textbf{1. Finding and using an inverse.}

Let
\[
f(x)=4x-7.
\]
\begin{enumerate}[label=(\alph*),leftmargin=0.35in]
\item Find $f^{-1}(x)$.
\answerblank{0.62in}
\item Find $f^{-1}(13)$.
\answerblank{0.42in}
\item Verify your inverse by simplifying $f\bigl(f^{-1}(x)\bigr)$.
\answerblank{0.68in}
\end{enumerate}

\textbf{2. Restricting a domain.}

Explain why $g(x)=x^2+2$ does not have an inverse function on all real numbers. Restrict the domain to $x\ge0$, find $g^{-1}(x)$, and state the domain of the inverse.
\answerblank{1.55in}
\newpage

% MODULE 5 CHECKIN PAGE 2
\assessmenttitle{5}{Inverses and Radicals -- Weekly Check-In, continued}
\studentline
\textbf{3. Solving a radical equation.}

Solve
\[
\sqrt{2x+3}=x.
\]
Show the algebra that produces all candidates, then check each candidate in the \emph{original} equation and identify any extraneous solution.
\answerblank{4.35in}

\end{document}
